\documentclass[10pt,a4paper]{amsart}
\usepackage[margin=19mm]{geometry}
\usepackage{amsmath,amssymb,mathtools}
\usepackage[scr=rsfs,cal=euler]{mathalfa}
\usepackage{xfrac}
\usepackage[unicode,hidelinks,bookmarks=false]{hyperref}
\newcommand{\ie}{i.e.\ }
\newcommand{\cf}{cf.\ }
\newcommand{\resp}{respectively\ }
\newcommand{\Subsection}{\subsection}
\providecommand{\nobreakdash}{\nobreak\hbox{-}}
\setlength{\emergencystretch}{2em}
\multlinegap0pt
\usepackage{amssymb}
\usepackage{xcolor}
\usepackage{bm}
\usepackage{tcolorbox}
\usepackage{mathtools}
\usepackage{tikz}
\newtheorem{lemma}{Lemma}
\title{Computing and Bounding the Number of Eulerian Orientations for Certain Classes of $4$-Regular Graphs}
\author{Evangelos Bartzos, Michalis Samaris}
\date{Source: arXiv:2609.06701v1 (CC BY 4.0)}
\begin{document}
\maketitle

\noindent\emph{Source and licence.} Computing and Bounding the Number of Eulerian Orientations for Certain Classes of $4$-Regular Graphs. Version arXiv:2609.06701v1 (CC BY 4.0), distributed by arXiv under the Creative Commons Attribution 4.0 International licence, http://creativecommons.org/licenses/by/4.0/, as recorded in the arXiv licence metadata for this exact version and shown on the abstract page https://arxiv.org/abs/2609.06701v1. Full licence notice: \texttt{licenses/arxiv-2609.06701-CC-BY-4.0.txt}.\par\smallskip

\noindent\textit{Editorial scope.} Counted unit: the article's lemma lem:count\_c (source0.tex lines 2862--2871). The article's complete original proof of it is delivered with it (source0.tex lines 2873--2911, one \texttt{proof} environment of 39 lines). No mathematics is written, repaired or re-derived: the statement and the proof are the article's own lines, in the article's own notation, and no retained line is substituted or re-flowed.  Retained uncounted as the source-local context the proof needs: lines 2818--2820.  Nothing was omitted from any retained span, so the package carries no omission record.  No retained line contains a citation, so the wrapper carries no bibliography and no bibliography entry.  No retained line contains a figure environment, an image-inclusion command or drawing code, so no image is delivered and the package declares no asset dependency.  Editorial formatting only, in addition: an AMS \texttt{amsart} preamble that restates the article's own declarations and macros used by the retained lines, namely the environments \texttt{lemma} on separate counters, together with the standard packages those lines need.  The retained environments print in this document as Lemma 1 in order of appearance, whereas the article prints the same items with the numbering of its own full document.  The complete native source of the article as distributed in the arXiv e-print for this exact version is bundled unchanged as \texttt{sources/209561/source0.tex}.
        \begin{lemma}\label{lem:structure 4regular}
            Each biconnected component of the induced subgraph $G[V_c]$ is a cycle.
        \end{lemma}

        \begin{lemma}\label{lem:count_c}
            Let $\mathcal{T}$ be a connected component of $G[V_c]$ which is not a single vertex.
            Let $V^2_{\mathcal{T}}$ and $V^4_{\mathcal{T}}$ be the subsets of vertices from $V^2_{c}$ and $V^4_{c}$ belonging to $V(\mathcal{T})$, respectively, and let $s$ be the number of cycles in $\mathcal{T}$.
            Then:
            \begin{itemize}
                \item[(i)]  $s=|V^4_{\mathcal{T}}|+1$.
                \item[(ii)] If every cycle in $\mathcal{T}$ is simple, then $|V^4_{\mathcal{T}}|\leq |V^2_{\mathcal{T}}|-3$.
                Equality holds when every cycle has length $3$.
            \end{itemize}
        \end{lemma}

        \begin{proof}
            We prove both (i) and (ii) by induction on the number of degree-$4$ vertices in a graph $\mathcal{T}$ where every vertex has degree $2$ or $4$.
            If $|V^4_{\mathcal{T}}|=0$ then $\mathcal{T}$ is a cycle. 
            Clearly, in $\mathcal{T}$ it holds that $s=|V^4_{\mathcal{T}}|+1=1$, satisfying (i).
            Furthermore, if this cycle is simple then it has at least 3 vertices.
            Thus, $|V^4_{\mathcal{T}}|\leq |V^2_{\mathcal{T}}|-3$. 
            The equality holds if the length of the cycle is 3, satisfying (ii).
            
        
            Now, assume that (i) and (ii) hold in  every graph $\mathcal{T}$ where every vertex has degree 2 or 4 and for the number of vertices of degree 4 $V^4_{\mathcal{T}}$  are strictly less than $k$ for some $k \geq 1$. 
            We will prove that they also hold when the graph has exactly $k$ vertices of degree $4$. 
            
            Let $v\in V^4_{\mathcal{T}}$, and let the subgraphs $\mathcal{T}_1\cup\{v\}$ and $\mathcal{T}_2\cup\{v\}$ be constructed as in the proof of Lemma~\ref{lem:structure 4regular}.
            Clearly, $\mathcal{T}_1\cup\{v\}$ and $\mathcal{T}_2\cup\{v\}$ satisfy the conditions of the induction hypothesis since, in each of them, every vertex has degree 2 or 4 and there are at most $k-1$ vertices degree 4.
            Let $V^2_{\mathcal{T}_1}$ and $V^4_{\mathcal{T}_1}$ be the vertices of degree 2 and 4 respectively in $\mathcal{T}_1\cup\{v\}$. 
            Similarly, let $V^2_{\mathcal{T}_2}$ and $V^4_{\mathcal{T}_2}$ be the vertices of degree 2 and 4 respectively in $\mathcal{T}_2\cup\{v\}$.
            
            Except for $v$, every vertex belonging in $\mathcal{T}_1\cup\{v\}$ or $\mathcal{T}_2\cup\{v\}$ has the same degree in the respective graph that belongs as in $\mathcal{T}$, whereas $v$ has degree 2 in both of them. 
            Let also  $s_1$ and $s_2$ be the number of cycles in $\mathcal{T}_1\cup\{v\}$ or $\mathcal{T}_2\cup\{v\}$ respectively. 
            Then, the following equalities hold:
        
        
            \begin{align*}
                |V^2_{\mathcal{T}}| &= (|V^2_{\mathcal{T}_1}|-1)+(|V^2_{\mathcal{T}_2}|-1) = |V^2_{\mathcal{T}_1}|+|V^2_{\mathcal{T}_2}|-2,  \\
                |V^4_{\mathcal{T}}| &= |V^4_{\mathcal{T}_1}|+|V^4_{\mathcal{T}_2}|+1, \text{ and } \\
                s &= s_1+s_2.
            \end{align*}
        
            \textbf{Proof of (i).} By the induction hypothesis, it holds that $s_1=|V^4_{\mathcal{T}_1}|+1$ and $s_2=|V^4_{\mathcal{T}_2}|+1$.
            Combining this with the above equalities yields:
            $$s = s_1+s_2 = (|V^4_{\mathcal{T}_1}|+1) + (|V^4_{\mathcal{T}_2}|+1) = |V^4_{\mathcal{T}}|+1.$$
        
            \textbf{Proof of (ii).} By the induction hypothesis, it holds that $|V^4_{\mathcal{T}_1}|\leq |V^2_{\mathcal{T}_1}|-3$ and $|V^4_{\mathcal{T}_2}|\leq |V^2_{\mathcal{T}_2}|-3$. 
            Combining this with the above equalities yields:
             $$|V^4_{\mathcal{T}}| = |V^4_{\mathcal{T}_1}|+|V^4_{\mathcal{T}_2}|+1 \leq (|V^2_{\mathcal{T}_1}|-3) + (|V^2_{\mathcal{T}_2}|-3) + 1 = |V^2_{\mathcal{T}}|-3.$$
        
            If $|V^4_{\mathcal{T}_1}|= |V^2_{\mathcal{T}_1}|-3$ and $|V^4_{\mathcal{T}_2}|= |V^2_{\mathcal{T}_2}|-3$, then every cycle in $\mathcal{T}_1$ and $\mathcal{T}_2$ has length $3$ by the induction hypothesis. 
            In this case, clearly every cycle in $\mathcal{T}$ also has length $3$, and the above inequality holds with equality, i.e., $|V^4_{\mathcal{T}}|=|V^2_{\mathcal{T}}|-3$.
        \end{proof}

\end{document}
